% chapters/ch2-intro-sec1.tex -- Introduction; 1. The Dynamical Equations, 1.1-1.4 (PDF 93-110)
\unnumberedsection{Introduction}\label{ch2:sec:intro}

Model rocketeers have been familiar with the concept of aerodynamic stability
for quite some time. As early as 1958 G. Harry Stine had published simplified
treatments of aerodynamic stability as it pertains to model rocket design, in
which the hobbyist was first made aware of the existence of a center of
pressure (\CP) and of the all-important relationship between it and the model's
center of gravity (\CG). In these early treatments the center of pressure was
approximated by the center of lateral area, and for the next eight years the
``cardboard cutout'' method remained the standard model builder's technique for
estimating the location of the \CP

The next major advance in the field came in 1966 at the Eighth National
Association of Rocketry Annual Meet, where James Barrowman of the National
Aeronautics and Space Administration's Sounding Rocket Branch unveiled an
algebraic method based on the theory of potential flow capable of determining
the center of pressure of a model rocket flying subsonically and at small
angles of attack to a high order of accuracy. Barrowman showed that the actual
\CP{} of a model rocket lies some distance aft of the location predicted by the
cutout method, and that therefore model rocketeers had been designing their
vehicles too conservatively.

All consideration of model rocket stability had thus far been confined
exclusively to its static manifestations, with the nature of the \CP-\CG{}
relationship being used to determine whether a rocket, once deflected from
facing directly into the relative airstream, would experience a moment tending
to return it to the undeflected state (in which case it would be considered
stable) or one tending to further deflect it (in which case it would be
considered unstable). Little attention had been paid to the details of the
process by which a stable rocket, once disturbed in its flight, restores itself
to alignment with its intended flight path, or to the process by which an
unstable rocket goes head over heels (unless, strangely enough, it is spinning
rapidly enough about its centerline). The \emph{statics} of stability had been
treated in admirable detail, but its \emph{dynamics} remained virtually
ignored.

This is not to say that our hobby had had \emph{no} exposure to stability
considered in a dynamic context. Luther W. Gurkin of the National Aeronautics
and Space Administration's Wallops Station had presented an excellent short
treatment of rocket dynamics to the contestants at the Sixth National
Association of Rocketry Annual Meet in 1964. Although copies of Gurkin's
``Basic Missile Aerodynamic Stability'' were distributed to a number of
interested persons at the meet, however, little was done to apply his results
specifically to model rockets for nearly four years. Our vehicles continued to
be subject to puzzling anomalies of behavior that could not be satisfactorily
explained by considerations of static stability alone. Sometimes they
``weathercocked''---flew directly into the prevailing wind at launch. In other
cases they would mysteriously tip to some random orientation at launch and
subsequently fly straight and true. Some models would oscillate excessively on
the way up; others hardly at all, and still others would experience a sort of
``semi-instability'' in which the nose circled violently about the intended
axis of flight.

This writer began to investigate the problem of dynamic stability early in
1968, working from a general consideration of the dynamical equations governing
the rotational motions of a streamlined projectile about its center of gravity
during free flight. Although the mathematical details of such an approach are
sometimes formidable, I felt that the greatest amount of information could only
be obtained from the most general analysis. In this section I am going to
endeavor to present the results of this analysis and to apply them to model
rockets by the use of the Barrowman equations, concluding with some suggestions
which should help the model rocketeer formulate designs which will both fly
stably and exhibit favorable dynamic behavior.

The dynamical equations describing the behavior we are interested in are
necessarily of a type called \emph{differential} equations, and as such require
the techniques of calculus for their solution. I want to emphasize as strongly
as I can, however, that it is \emph{not necessary} that the reader understand
calculus in order to follow the presentation. To the engineer, calculus is
fundamentally a tool that enables him to obtain \emph{algebraic} equations
describing the behavior of the system he is investigating. Care has therefore
been taken to emphasize the \emph{algebraic results} of calculus-dependent
derivations, with the calculus operations being considered as formulae for
altering differential equations to algebraic equations. No reader who has had
his second year of high school algebra, and perhaps some exposure to analytic
geometry, should have any difficulty in understanding the text.

As the exact forms of the equations governing the rotational motions of
projectiles subject to aerodynamic moments are quite complex and introduce
fundamental mathematical barriers to our obtaining physical solutions, I have
made a number of approximations in order to cast the equations into a more
readily soluble form. Approximations of the kind made here are quite common in
the solution of the differential equations encountered in all branches of
mathematical physics. Called \emph{linearizations}, they involve physical
and/or geometrical reasoning by which the investigator can neglect or modify
certain of the characteristics of his mathematics in ways not immediately
determinable or obvious from the mathematics itself. I have taken care to
identify each such approximation and to give its physical justification in
order to keep the treatment as basic as possible.

Certain portions of this analysis parallel Luther Gurkin's ``Basic Missile
Aerodynamic Stability'', while others use information contained in James
Barrowman's ``Calculating the Center of Pressure of a Rocket'', a National
Aeronautics and Space Administration Educational Services Office pamphlet
embodying the algebraic results of his analysis. In order to facilitate
correlation of this treatment with theirs, appropriate references have been
included.

\section{The Dynamical Equations}\label{ch2:sec:1}

\subsection{Euler's Angles}\label{ch2:sec:1.1}

Suppose we have a rocket which has been rotated about some set of mutually
perpendicular axes fixed in space: A, B, and C. We can speak of this rotational
displacement in a quantitative way if we consider a second set of axes, D, E,
and F, to have been fixed in the \emph{rocket}, with origin at the rocket's
\CG, and with directions coincident with A,B, and C before the rotation began,
and to have \emph{remained} fixed in the rocket as it rotated. We can always
uniquely determine the final orientation of the rocket if we agree to abide by
the following rule: that, in undergoing any given rotation, the rocket first
\emph{yaws} through an angle $\alpha_D$ about axis D; then \emph{pitches}
through an angle $\alpha_E$ about axis E; and finally \emph{rolls} through an
angle $\alpha_F$ about axis F. This process is illustrated in
Figure~\ref{ch2:fig:1}.

\begin{figure}[htbp]
  \centering
  \includegraphics{figures/v2/ch2/fig01.pdf}
  \caption{A set of Euler's angles for specifying the rotational position of a
  model rocket. $\alpha_D$ is the angle of \emph{yaw}, $\alpha_E$ is the angle
  of \emph{pitch}, and $\alpha_F$ is the angle of \emph{roll}. The origin of
  the coordinate system is taken as the center of gravity of the rocket.}
  \label{ch2:fig:1}
\end{figure}

The \emph{body axes} D,E, and F first rotate about D through angle $\alpha_D$.
Throughout this first rotation D coincides with axis A of the \emph{space axes}
and E and F remain in the plane defined by B and C. At its completion axis F
coincides with the dashed line \emph{Of} and axis E with \emph{Oe}. Next the
body axes rotate about the \emph{new position} of E (that is, about \emph{Oe})
through an angle $\alpha_E$. This rotation occurs in the plane of A and
\emph{Of}, and at its conclusion axis F is in its final position as shown and
axis D lies along line \emph{Od}. E, of course, is still along \emph{Oe}.
Finally the body axes execute the roll through angle $\alpha_F$ about F, in the
plane defined by \emph{Oe} and \emph{Od}, bringing all the axes into their
final positions.

It is important that this order of rotations be observed; if the order in which
the yaw, pitch, and roll occur is changed the final position of the rocket will
be different. A mathematician would say that angular displacements, or
rotations, are \emph{not vector quantities} because, although they specify both
a magnitude and a direction, they do not \emph{commute in addition}. The three
``locating angles'' $\alpha_D$, $\alpha_E$, and $\alpha_F$ for specifying the
rotational displacement of a solid body are called a \emph{set of Euler's
angles} after their discoverer, the Swiss mathematician and physicist Leonhard
Euler (1707--1783). You can see that if for any reason we wish to postulate an
\emph{intermediate} set of axes which will perform some, but not all, of the
movements of the rocket, we can find the rocket's final position by determining
\emph{first} the position of the intermediate axes with respect to the space
axes, and \emph{then} the position of the body axes with respect to the
intermediate axes. At this point such a procedure seems a meaningless
complication, but it will be very useful later on.

\subsection{Angular Velocity}\label{ch2:sec:1.2}

If you imagine the Euler's angles of a certain rotation becoming \emph{very
small}, so that the rocket is barely turned from its original position, you may
notice a curious and very useful fact: the order of the angular displacements
is no longer of such great importance. If, when the angles were large, the
rocket was considered to first pitch, then yaw, and finally roll or to perform
the rotations in \emph{any} order other than the prescribed one of yaw, pitch,
roll, we obtained a different final orientation in each case. Now, however, the
effect of such alterations in the order of rotations is very slight; nearly the
same result is obtained no matter what the order in which the rotations occur.
If the rotations are allowed to become \emph{infinitesimally} small, so that
they are, in the terminology of calculus, \emph{differential quantities}, the
equality becomes exact. Thus, differential rotations \emph{are} vector
quantities; they possess both a magnitude (though slight, to be sure) and a
direction, and they \emph{do} commute in addition. Differentials are denoted by
a lower-case letter $d$ in front of the quantity in question; differential
angular displacements are thus $d\alpha_D$, $d\alpha_E$, and $d\alpha_F$.

Now suppose the rocket is turning continuously, so that its Euler's angles keep
changing as time goes on. This can be represented mathematically by stating
that, in every \emph{differential element} of time $dt$, the rocket experiences
differential angular displacements $d\alpha_D$, $d\alpha_E$, and $d\alpha_F$.
We then form the fractions
\begin{equation}\label{ch2:eq:1}
\begin{split}
\omega_D &= \frac{d\alpha_D}{dt}\\
\omega_E &= \frac{d\alpha_E}{dt}\\
\omega_F &= \frac{d\alpha_F}{dt}
\end{split}
\end{equation}
and define
\begin{itemize}[nosep,leftmargin=*,label={}]
  \item $\omega_D \equiv$ angular velocity about D, or D-component of angular
  velocity
  \item $\omega_E \equiv$ angular velocity about E, or E-component of angular
  velocity
  \item $\omega_F \equiv$ angular velocity about F, or F-component of angular
  velocity
\end{itemize}
Angular velocity is called the \emph{derivative with respect to time} of
angular displacement, or the \emph{time rate of change} of angular
displacement. Angular velocities are vector quantities: they have both
magnitude and sign, they do commute in addition, and all vector operations such
as cross product, dot product, and coordinate transformation apply to them.
Specifically, we represent them in a so-called \emph{right-handed} coordinate
system as \emph{positive} about a given axis when they cause a rotation such
that a screw with right-handed threads would advance along that axis in the
positive direction, \emph{negative} when such a progression would be in the
negative direction. A right-handed system itself, such as our A,B,C axis set,
is established by the requirement that the right-handed screw turned about its
axis from A toward B advance positively along C. This convention is illustrated
in Figure~\ref{ch2:fig:2}, while Figure~\ref{ch2:fig:3} shows a rocket
undergoing a rotation in which all components of angular velocity are positive.
An angular velocity is conveniently represented by an arrow along the axis
about which it occurs, of direction determined by the sign of the component,
and of length proportional to its magnitude. This allows all the common vector
operations to be performed. Figure~\ref{ch2:fig:4} shows this convention in
operation.

\begin{figure}[htbp]
  \centering
  \includegraphics{figures/v2/ch2/fig02.pdf}
  \caption{Definition of a right-handed coordinate system. Turning a screw with
  right-hand threads as if axis A were being turned toward axis B causes the
  screw to advance positively along axis C.}
  \label{ch2:fig:2}
\end{figure}

\begin{figure}[htbp]
  \centering
  \includegraphics{figures/v2/ch2/fig03.pdf}
  \caption{Positive components of angular velocity. $\omega_D$ is a positive
  \emph{yaw rate}, $\omega_E$ is a positive \emph{pitch rate}, and $\omega_F$
  is a positive \emph{roll rate}.}
  \label{ch2:fig:3}
\end{figure}

\begin{figure}[htbp]
  \centering
  \includegraphics{figures/v2/ch2/fig04.pdf}
  \caption{Vector representation of positive angular velocity components.
  Compare with Figures~\ref{ch2:fig:2} and~\ref{ch2:fig:3} to see how this
  representation is suggested by the direction of advance of a right-handed
  screw.}
  \label{ch2:fig:4}
\end{figure}

It is now necessary to go back and pick up a few loose ends in order to
introduce some concepts which will be useful later on. First, there is the
matter of defining positive and negative rotations; this must be done to obtain
physically meaningful results from any investigation of rocket motion even
though angular displacements do not have all the properties required of
vectors. This is accomplished directly by the analogy of the right-handed screw
through replacing the ``turning direction'' of the screw with the angular
displacement. Secondly, the derivative relationship of angular velocity to
angular displacement needs a little more explanation. While physical reasoning
using differential quantities leads directly to the general form of
equations~\eqref{ch2:eq:1}, this technique is not very useful when the question
is: ``given the angular displacement as a function of time (that is, in an
analytical formula), compute the angular velocity''. In this case the relation
\begin{equation*}
\omega = \frac{d\alpha}{dt}
\end{equation*}
is taken to mean: ``given $\alpha(t)$ (the meaning of this notation is
``$\alpha$ as a function of time'' and it is read ``$\alpha$ of $t$''), apply a
known rule, or formula, for \emph{differentiating it} with respect to time and
obtain $\omega(t)$''. The expression, or formula, for $\alpha(t)$ is subject to
certain restrictions for this method to work properly, but these need not
concern you in anything discussed here. I am going to \emph{list the formulae}
for the derivatives needed in this treatment wherever they appear, so that it
will not be necessary to know them in order to follow the discussion.

\subsection{Applied Moments, Angular Accelerations, and Moment of Inertia}\label{ch2:sec:1.3}

Moment, or torque, is to rotational motion as force is to linear motion. It is
the cause of all changes in the state of the rotational motion of a physical
body. The simplest kind of relation between an applied moment and the resulting
angular motion is identical in form to Newton's second law of motion for
translational displacement: $\vec{F} = m\vec{a}$. To see this consider a
flywheel, initially at rest, mounted in an axle held in frictionless bearings.
At some time arbitrarily designated as $t = 0$ we begin to apply a constant
moment about the shaft. An angular velocity will then arise which starts from
zero and increases linearly (i.e., at a constant rate) with time. The rate is
directly proportional to the applied moment and inversely proportional to a
property of the mass distribution in the wheel called its \emph{moment of
inertia} about its axis of radial symmetry. The relationship between an applied
moment and the resulting \emph{angular acceleration} is written
\begin{equation}\label{ch2:eq:2}
M = I\gamma
\end{equation}
where
\begin{itemize}[nosep,leftmargin=*,label={}]
  \item $M$ = applied moment
  \item $I$ = moment of inertia
  \item $\gamma$ = angular acceleration
\end{itemize}
Note that the angular acceleration $\gamma$ is the time rate of change of the
angular velocity $\omega$, just as $\omega$ is the time rate of change of the
angular displacement $\alpha$. This sequence of derivative relationships is
written
\begin{equation}\label{ch2:eq:3}
\begin{split}
\gamma &= \frac{d\omega}{dt}\\
&= \frac{d}{dt}\left(\frac{d\alpha}{dt}\right)\\
&= \frac{d^{2}\alpha}{dt^{2}}
\end{split}
\end{equation}
where the expression $d^{2}\alpha/dt^{2}$ is read, ``the second derivative of
$\alpha$ with respect to $t$''. From a practical standpoint it means that the
rule for differentiating $\alpha(t)$ has been applied twice in succession in
order to obtain $\gamma(t)$. Thus $\gamma$, $\omega$, and $\alpha$ are all
related by derivatives with respect to time, a property of great value to
analysis. In the case of a constant applied moment, for instance, the angular
acceleration is given by
\begin{equation*}
\gamma = \frac{M}{I}
\end{equation*}
while any mathematician could tell you that the angular velocity is
\begin{equation*}
\omega = \frac{M}{I}\,t
\end{equation*}
and the angular displacement after a time $t$ is
\begin{equation*}
\alpha = \frac{M}{2I}\,t^{2}
\end{equation*}
These relations are illustrated in Figure~\ref{ch2:fig:5}. Note also that both
moment and angular acceleration are vectors, while moment of inertia, like
mass, is a scalar. Moment of inertia gets its name not from the fact that an
object exerts an equal and opposite moment on the cause of its angular
acceleration (although it does), but from the fact that it is, mathematically
speaking, the ``second moment'' of mass about a given axis, a term arising from
the calculus formula for $I$. This analysis will not present the \emph{integral
form} that must be used to calculate $I$ for the most general case, but the
results for a few particular objects will appear later in the treatment.

\begin{figure}[htbp]
  \centering
  \includegraphics{figures/v2/ch2/fig05.pdf}
  \caption{Angular acceleration of a flywheel. At time $t = 0$ a constant
  moment $M$ is applied to a flywheel whose moment of inertia about the axis is
  $I$, producing an angular acceleration $\gamma = M/I$.}
  \label{ch2:fig:5}
\end{figure}

\subsection{Euler's Dynamical Equations}\label{ch2:sec:1.4}

The general three-dimensional angular motion of rigid bodies in response to
general applied moments is far more complicated than the simple flywheel
example given above. Their derivation is of no interest to model rocketeers,
and I have therefore omitted it entirely from this presentation. The results of
interest to us are the dynamical equations for a body which has sufficient
symmetry for the directions of the so-called \emph{principal axes} to be
geometrically obvious, and which is free to rotate about its center of mass in
any direction. These equations, for a body with principal body axes D,E,F,
rotating with respect to axes A,B,C fixed in space, are
\begin{equation}\label{ch2:eq:4}
\begin{split}
M_D &= I_D\frac{d\omega_D}{dt} - (I_E - I_F)\,\omega_E\omega_F\\
M_E &= I_E\frac{d\omega_E}{dt} - (I_F - I_D)\,\omega_F\omega_D\\
M_F &= I_F\frac{d\omega_F}{dt} - (I_D - I_E)\,\omega_D\omega_E
\end{split}
\end{equation}
where $M_D$, $M_E$, and $M_F$ are moments about D,E, and F impressed by
external agencies and $I_D$, $I_E$, and $I_F$ are the moments of inertia taken
about axes D,E, and F respectively. These equations, also due to Euler, were a
landmark in the history of classical dynamics and are named Euler's dynamical
equations in his honor. For bodies having trigonal or greater mass symmetry
about the longitudinal axis (this category includes most rockets) we have
\begin{align*}
I_D = I_E &\equiv I_L\\
I_F &\equiv I_R
\end{align*}
and the equations reduce to
\begin{equation}\label{ch2:eq:5}
\begin{split}
M_D &= I_L\frac{d\omega_D}{dt} - (I_L - I_R)\,\omega_E\omega_F\\
M_E &= I_L\frac{d\omega_E}{dt} - (I_R - I_L)\,\omega_F\omega_D\\
M_F &= I_R\frac{d\omega_F}{dt}
\end{split}
\end{equation}
It will become clear in a little while that these forms of the Euler equations
are a bit inconvenient to use. It will be preferable to construct an
intermediate system of axes X, Y, and Z which follow the rocket in yaw and
pitch, but do not roll. Thus, Z always coincides with F but X and Y do not
coincide with D and E unless $\alpha_F = 0$. Therefore,
\begin{align*}
\alpha_X &= \alpha_D\\
\alpha_Y &= \alpha_E\\
\alpha_Z &= 0 \ne \alpha_F
\end{align*}
This set of axes is shown in Figure~\ref{ch2:fig:6}. Since the rocket is
symmetrical about Z, the intermediate axes remain principal regardless of the
roll angle $\alpha_F$. Let the angular velocity components of the intermediate
axes be denoted by $\Omega_X$, $\Omega_Y$, and $\Omega_Z$. Then we have
\begin{align*}
\Omega_X &= \omega_D\\
\Omega_Y &= \omega_E\\
\Omega_Z &= 0 \ne \omega_F
\end{align*}
Since the Z axis always coincides with the F axis, the notation can be
simplified to include only three coordinate variables by writing
\begin{equation*}
\omega_Z \equiv \omega_F
\end{equation*}
but we must be careful not to confuse $\omega_Z$ (the roll rate of the rocket)
with $\Omega_Z$ (the roll rate of the intermediate coordinates, which is kept
zero). The dynamical equations now become, for a body with trigonal or greater
mass symmetry about Z,
\begin{equation}\label{ch2:eq:6}
\begin{split}
M_X &= I_L\frac{d\Omega_X}{dt} + I_R\,\Omega_Y\omega_Z\\
M_Y &= I_L\frac{d\Omega_Y}{dt} - I_R\,\Omega_X\omega_Z\\
M_Z &= I_R\frac{d\omega_Z}{dt}
\end{split}
\end{equation}

\begin{figure}[htbp]
  \centering
  \includegraphics{figures/v2/ch2/fig06.pdf}
  \caption{Space axes, intermediate axes, and body axes. The space axes (A,B,C)
  are rotationally fixed with respect to inertial space. The body axes (D,E,F)
  are fixed in the airframe of the rocket and follow it in yaw, pitch, and
  roll, while the intermediate axes (X,Y,Z) follow the rocket in yaw and pitch
  but maintain a zero roll angle.}
  \label{ch2:fig:6}
\end{figure}
